\documentclass[11pt]{article}
\usepackage[margin=25mm]{geometry}
\usepackage{fontspec}
\setmainfont{Latin Modern Roman}
\usepackage{amsmath,amssymb,mathtools}
\usepackage{unicode-math}
\setmathfont{Latin Modern Math}
\usepackage{xurl,hyperref}
\hypersetup{colorlinks=true,urlcolor=blue}
\setcounter{secnumdepth}{0}
\setlength{\parindent}{0pt}
\setlength{\parskip}{5pt}
\setlength{\emergencystretch}{3em}
\newcommand{\sref}[2]{\hyperlink{source:#1:#2}{[S#2]}}
\newcommand{\cataloguescope}[1]{\paragraph{Recorded target.}#1}
\newcommand{\R}{\mathbb R}
\newcommand{\C}{\mathbb C}
\newcommand{\N}{\mathbb N}
\newcommand{\Z}{\mathbb Z}
\newcommand{\Q}{\mathbb Q}
\newcommand{\F}{\mathbb F}
\newcommand{\eps}{\varepsilon}
\begin{document}
\section{When constrained minimizers gain symmetry}\label{statement}

\subsection{Definitions and mathematical statement}

\paragraph{Shared notation.}$\mathbb N=\{1,2,\ldots\}$; the standard number fields have their usual meanings. All additional conventions are stated locally.


\paragraph{Definitions and notation.}In the Newtonian $n$-body problem, positions $q_i\in\mathbb R^d$ and masses $m_i>0$ obey $m_i\ddot q_i=\nabla_{q_i}U$, where $U(q)=\sum_{i<j}m_im_j/|q_i-q_j|$ and the gravitational constant is normalized to one. Here $T=\tfrac12\sum_i m_i|\dot q_i|^2$, $E=T-U$, and a collision has $q_i=q_j$. The center of mass can be fixed at zero; moment of inertia is $I=\sum_i m_i|q_i|^2$. For a fixed period $P$, the action is $\mathcal A(q)=\int_0^P(T+U)\,dt$. A collision-free minimizer in an admissible loop class satisfies the Newtonian equations when variations are permitted. The period, constraints and closure of the class must be specified: a minimizing sequence can escape or approach collisions. A symmetry class consists of loops invariant under a specified combined action on time, space and labels. A loop’s full stabilizer may contain this prescribed group strictly. Extra symmetry means this strict enlargement; it need not follow merely from existence in the admissible class.

\paragraph{Catalogue wording retained for traceability.}
Problem 9. Often the minimizers in a symmetry class has more symmetry that the class asks for. Is that a coincidence? (Alain Chenciner)
\paragraph{Scope and source qualification.}The source asks for an explanation of a phenomenon and gives examples, without a universal quantifier or precise class of groups. It must not be rewritten as every minimizer has strictly more symmetry. The catalogue title/status was paired with a machine-generated research attempt; that attempt is not peer-reviewed evidence of settlement. The statement covers the original source questions rather than only the title’s subcase.
\subsection{Short English statement}
Explain and characterize why constrained action minimizers often have extra symmetries.
\subsection{Research attempt}
\paragraph{Model, date and investigated scope.}
OpenAI GPT-6 Astra (configured model \texttt{gpt-6-astra}), ultra reasoning; generated 11 October 2026. Proved the uniqueness-based symmetry mechanism and supplied counterexamples to symmetry-only and naive convex-averaging arguments. No n-body class-specific extra-symmetry theorem is proved.

\paragraph{Sources and literature check.}
Read the supplied complete definition, catalogue wording, scope qualifications and source list for this canonical ID. The original-source citations are inherited as supplied references; no fresh primary-source or novelty search was performed for this attempt. No current open-status claim is inferred from this calculation.

\paragraph{A mechanism from uniqueness.}Let a group \(G\) preserve a functional \(\mathcal A\) and an admissible set \(X\). If \(q\) is its unique minimizer in \(X\), then \(gq=q\) for every \(g\in G\): invariance makes \(gq\) another minimizer, so uniqueness forces equality. Applied to a class fixed by a smaller group \(H\), this gives extra symmetry when a larger normalizing group preserves that class and the minimizer is unique. If minimizers are unique only modulo time translation or rotation, the conclusion is invariance of that orbit, and an actual pointwise stabilizer requires choosing and checking representatives.
\paragraph{Why symmetry alone does not suffice.}The smooth even function \(f(x)=(x^2-1)^2\) is invariant under \(x\mapsto-x\), but its two minimizers \(\pm1\) are exchanged and neither is fixed. Moreover averaging symmetric loops cannot generally rescue the argument for Newtonian action: the Hessian of \(|x|^{-1}\) is \(3xx^T/|x|^5-I/|x|^3\), with negative tangential eigenvalues. Thus the potential is not convex, and averaging can increase action or even create collisions.
\paragraph{First gap.}For a concrete n-body symmetry class one must prove uniqueness (possibly modulo continuous symmetries), or identify a different inequality forcing enlargement of the stabilizer. Neither is established for the source's examples. The elementary uniqueness mechanism explains one possible reason for extra symmetry and the counterexample shows why no universal claim follows from invariance alone. The source asks about a phenomenon, not the false theorem that every constrained minimizer has strictly greater symmetry.

\paragraph{Verification and adversarial audit.}
Checked group preservation of the class, uniqueness versus orbit uniqueness, two-minimum counterexample and the Newtonian Hessian signs. This was an LLM adversarial audit, not independent human verification; no theorem prover was used. Prompts A then B applied to this ID: Checked group preservation of the class, uniqueness versus orbit uniqueness, two-minimum counterexample and the Newtonian Hessian signs. Scope remains limited: Establish uniqueness modulo the correct continuous symmetries or a separate action inequality for an actual n-body admissible class.

Coordinator checked the uniqueness symmetry argument, exchanged-minima counterexample and Newtonian potential Hessian signs. This additional check was model-assisted, not independent human verification.

\paragraph{Tools, process and contributions.}
Python 3 and local shell used to read supplied files and serialize this manuscript; mathematics derived symbolically in the text, with no numerical experiment or formal proof assistant. Contribution prompt A was followed by prompt B. GPT-6 Astra with ultra reasoning drafted the research and performed its stated self-audit; any additional coordinator contribution is identified in the verification and audit. A GPT-6-based Codex coordinator assembled the manuscript and checked repository integration; its exact variant and reasoning configuration were not exposed. No human mathematical contribution, independent human proof check or formal proof certification is supplied. Complete derivations and the consulted source references are included here; private reasoning transcripts are not published.

\paragraph{Final status and first remaining gap.}
\textbf{UNRESOLVED}. Establish uniqueness modulo the correct continuous symmetries or a separate action inequality for an actual n-body admissible class.

\section{Completion Estimate}
\noindent\textbf{COMPLETION: 0\%}
The calculation is an elementary benchmark or conditional reduction; it establishes no claimed new portion of the complete source target.

\subsection{Sources}
\begin{itemize}
\item[\hypertarget{source:statement:1}{S1}] ULAM.AI and the UnsolvedMath Contributors, UnsolvedMath, catalogue entry. Catalogue curation under CC BY 4.0; definitions and English exposition added locally.
\url{https://huggingface.co/datasets/ulamai/UnsolvedMath}.
\item[\hypertarget{source:statement:2}{S2}] American Institute of Mathematics, Variational Methods in Celestial Mechanics; notes by Kuo-Chang Chen, version 22 June 2003, source item 9.
\url{https://aimath.org/WWN/varcelest/varcelest.pdf}.
\end{itemize}
\label{end:statement}
\end{document}
